\documentclass[aspectratio=169]{beamer}

\usetheme{Warsaw}
\usecolortheme{seahorse}
\setbeamertemplate{navigation symbols}{}
\setbeamertemplate{footline}[frame number]

\usepackage[T1]{fontenc}
\usepackage[utf8]{inputenc}
\usepackage{lmodern}
\usepackage{booktabs}
\usepackage{tabularx}
\usepackage{array}

\title{MLLM-SHAP Research Direction Plan}
\subtitle{Roadmap, Compute Demand, and Execution Timeline}
\author{Warsaw University of Technology}
\date{\today}

\begin{document}

\begin{frame}
  \titlepage
\end{frame}

\begin{frame}{Talk map (10 min)}
  \begin{itemize}
    \item Current plan status: blockers and execution gates
    \item Research tracks: audio/multimodal, text-to-text, extended program, infrastructure
    \item Compute demand model with explicit assumptions and dataset scales
    \item Integrated timeline with interleaved workstreams
    \item Scope for the master's thesis and conference papers
    \item Team extension
  \end{itemize}
\end{frame}

\begin{frame}{Research scope}
  \begin{itemize}
    \item \textbf{Reproducibility requirement:} versioned artifacts, pinned environment, benchmark protocol.
    \item \textbf{T2T track:} estimand correctness, exact-vs-approx estimator validation, utility-function ablations.
    \item \textbf{SGPA track:} segmentation quality, masking ablations, faithfulness, stability, human validation.
    \item \textbf{Cross-model requirement:} primary model + secondary local model + secondary API model.
    \item \textbf{Cross-lingual requirement:} English + at least two typologically distinct languages.
  \end{itemize}
\end{frame}

\begin{frame}{Current blockers and execution risks}
  \begin{enumerate}
    \item \textbf{Faithfulness evidence gap:} core faithfulness experiments are still incomplete and must be finished before strong claims.
    \item \textbf{Hyperparameter-grid validation:} segmentation/masking hyperparameters need full validation on natural speech and multi-annotator protocol.
    \item \textbf{Compute bottleneck:} planned experiment matrix is larger than single-GPU practical capacity; cluster scheduling is required.
    \item \textbf{Rerun risk:} if ablations change the best masking strategy, previously reported results must be recomputed.
  \end{enumerate}
\end{frame}

\begin{frame}{Research steps: text-to-text core}
  \begin{itemize}
    \item \textbf{Phase 0:} immediate data integrity audit for estimand identity (Shapley vs Banzhaf) across existing T2T outputs.
    \item \textbf{Phase 1--2:} freeze 3-model text matrix and held-out prompt suite; enforce connector and masking invariants.
    \item \textbf{Phase 3:} synthetic toy games with exact ground truth plus convergence checks against a precise explainer reference.
    \item \textbf{Phase 4:} utility-function/scoring ablations (greedy similarity, log-probability, distribution-aware alternatives).
    \item \textbf{Phase 5--6:} faithfulness and robustness tests (AOPC/AUC, rank stability, FDR correction), then benchmark/reporting hardening.
  \end{itemize}
\end{frame}

\begin{frame}{Research steps: infrastructure and tooling}
  \begin{itemize}
    \item \textbf{Data/artifacts:} move heavy results out of git, add artifact metadata (pipeline version, commit, config hash).
    \item \textbf{Framework outputs:} per-sample structured records with masks, responses, and hardware traces.
    \item \textbf{Performance:} optimize bottlenecks in masking/sampling and reduce RAM footprint for 10k+ coalition calls.
    \item \textbf{Reporting/visualization:} reproducible benchmark sidecars and notebook/GUI views for faithfulness and variance comparisons.
    \item \textbf{Role in program:} infrastructure track that enables reproducible scale-up of all scientific phases above.
  \end{itemize}
\end{frame}

\begin{frame}{Research steps: audio/multimodal core}
  \begin{itemize}
    \item \textbf{Phase 0 (gate):} lock model matrix, resolve Shapley/Banzhaf labeling, formalize multimodal coalition game.
    \item \textbf{Phase 1:} verify additivity assumptions, retune segmentation hyperparameters on real speech with multi-annotator protocol, add audio preprocessing.
    \item \textbf{Phase 2:} CTC alignment benchmark vs naive baselines; masking strategy ablation (deletion vs silence vs ambient noise) as go/no-go gate.
    \item \textbf{Phase 3:} full-dataset execution, cross-lingual runs, replication on secondary models, and power-analysis-driven sample sizing.
    \item \textbf{Phase 4--6:} faithfulness/stability/K-sensitivity, human validation, theory/reproducibility, and explicit failure-mode characterization.
  \end{itemize}
\end{frame}

\begin{frame}{Research steps: extended research program}
  \begin{itemize}
    \item \textbf{Phase 1:} multimodal interaction values (pairwise tractable SII-style studies), broader audio K-sensitivity, and plausibility/OOD metrics.
    \item \textbf{Phase 2:} generalization beyond current setup: encoder backbones, task families, gradient/attention baselines, model-size effects.
    \item \textbf{Phase 2:} modality trade-off analysis and fine-tuning impact on attribution rank structure.
    \item \textbf{Phase 3:} longer-horizon method work: learned coalition sampling, hierarchical grouping, KL-based scoring, contrastive explanations.
    \item \textbf{Phase 3:} input-space vs latent-space explanation layer comparison where internal representations are available.
  \end{itemize}
\end{frame}

\begin{frame}{Conference and thesis status}
  \begin{itemize}
    \item \textbf{INTERSPEECH 2025}
      \begin{itemize}
        \item SGPA method paper focused on spectrogram-guided phonetic alignment for feasible audio SV analysis.
        \item Core content: feasibility bottleneck diagnosis, 4-stage SGPA pipeline, segmentation diagnostics, and evaluation reductions (calls/time) for audio explainability.
        \item \textbf{Status:} rebuttal stage; likely rejection.
      \end{itemize}
    \item \textbf{ACL 2026 System Demonstration}
      \begin{itemize}
        \item \texttt{mllm-shap} platform/demo paper focused on system architecture and tooling.
        \item Core content: multimodal masking, multi-turn token-role tracking, estimator suite (Exact/MC/CC/Neyman/Hierarchical), GUI workflow, and package-level benchmarking.
        \item \textbf{Status:} accepted.
      \end{itemize}
    \item \textbf{Bachelor thesis}
      \begin{itemize}
        \item Thesis track aligned with the explainability platform and research program.
        \item \textbf{Status:} final grade 5.
      \end{itemize}
  \end{itemize}
\end{frame}

\begin{frame}{What is already evidenced in submitted work}
  \begin{itemize}
    \item \textbf{From the SGPA conference paper:} demonstrated feasibility improvement for audio explainability, including strong reduction in model calls per sample and improved boundary-quality diagnostics.
    \item \textbf{From the ACL system demo paper:} end-to-end platform scope is already demonstrated (multimodal masking, estimator suite, GUI workflow, package integration).
    \item \textbf{Reviewer perspective (current signal):} engineering feasibility and system integration are seen as strengths; faithfulness depth, broader validation, and stronger replication are the main pressure points.
    \item \textbf{Planning implication:} next steps are scale-up, replication, and stricter validation, rather than proof-of-concept work.
  \end{itemize}
\end{frame}

\begin{frame}{Dataset sizes driving compute volume}
  \scriptsize
  \begin{tabularx}{\textwidth}{>{\raggedright\arraybackslash}p{3.3cm} >{\raggedright\arraybackslash}X >{\raggedright\arraybackslash}p{2.0cm}}
    \toprule
    \textbf{Dataset block} & \textbf{Planned scope} & \textbf{Size} \\
    \midrule
    VoiceBench single-sentence subset & initial controlled English set with original + male TTS + female TTS variants & \(500\) \\
    LibriSpeech single-sentence subset & initial controlled natural-speech English set & \(500\) \\
    Full VoiceBench corpus & mandatory full run in large-scale phase & \(854\) \\
    Full single-sentence corpus & mandatory full run in large-scale phase & \(448\) \\
    Cross-lingual evaluation & at least two additional languages, minimum planned sample count & \(\ge 2 \times 100\) \\
    Power-analysis pilot set & pilot for effect-size estimation before full commitments & \(20\) \\
    Stability protocol set & repeated-seed and temperature protocols on fixed samples & \(20\) \\
    K-sensitivity validation set & \(K \in \{50,100,200,500,1000\}\) on fixed audio samples & \(20\) \\
    Human-study set & user validation study for highlighted segments & \(n \ge 15\) participants \\
    T2T held-out prompt suite & 6 prompt buckets (QA, extraction, rewrite, constrained gen, multi-turn, long-context) & \(50\)--\(100\) per bucket \\
    \bottomrule
  \end{tabularx}
\end{frame}

\begin{frame}{Lower-bound run count and runtime translation}
  \scriptsize
  \begin{tabularx}{\textwidth}{>{\raggedright\arraybackslash}p{3.0cm} >{\raggedright\arraybackslash}X >{\raggedright\arraybackslash}p{2.0cm}}
    \toprule
    \textbf{Component} & \textbf{Derivation from plan minima} & \textbf{Sample-runs} \\
    \midrule
    Core full-corpus faithfulness matrix & \(1302 \times 3\) methods \(\times 2\) full models & \(7812\) \\
    Cross-lingual minimum matrix & \((2 \times 100)\times 3\) methods \(\times 2\) full models & \(1200\) \\
    Stability + K checks & \((20\times10) + (20\times5)\) from repeated seeds and K-sweep & \(300\) \\
    Additional explicit subsets & 30-sample independence + 100-sample granularity + 40 failure-mode cases & \(170\) \\
    \midrule
    \textbf{Audio lower bound} & sum of explicit minima above (before secondary B and reruns) & \(\mathbf{9482}\) \\
    \bottomrule
  \end{tabularx}
\vspace{0.35em}
\begin{itemize}
  \item \textbf{Definition:} one \emph{sample-run} means one explained sample for one method-model configuration.
  \item \textbf{No-SGPA runtime mapping:} \(9482\) sample-runs \(\approx 9482\) GPU-hours \(\approx 395\) GPU-days (single GPU).
  \item \textbf{Mixed-runtime mapping:} if only the SGPA branch gets large speedup while native and LOO stay near baseline cost, total remains on the order of \(\sim 6.4\)k GPU-hours.
  \item \textbf{T2T suite scale note:} \(6\) buckets \(\times (50\text{ to }100)\) prompts gives \(300\) to \(600\) prompts; replicated across 2 research models and multiple explainers/scoring backends, this adds a substantial second compute block.
  \item \textbf{Lower-bound interpretation:} \(\sim 10\)k GPU-hours covers only the SGPA-focused lower bound; remaining components and the T2T grid add further demand.
\end{itemize}
\end{frame}

\begin{frame}{Timeline (interleaved by workstream)}
  \scriptsize
  \begin{tabularx}{\textwidth}{>{\raggedright\arraybackslash}p{2.2cm} >{\raggedright\arraybackslash}X}
    \toprule
    \textbf{Window} & \textbf{Primary focus and sequencing} \\
    \midrule
    Months 1--2 & \textbf{Fix package and unblock execution:} connector stability, experiment metadata/reporting, estimand labeling, baseline reproducibility checks. \\
    Months 3--5 & \textbf{T2T full grid:} held-out prompt-suite matrix, explainer/scoring ablations, two-model replication, significance/FDR-ready outputs. \\
    Months 5--8 & \textbf{SGPA core:} segmentation and masking ablations, alignment benchmarking, full-corpus faithfulness/stability/K-sensitivity runs. \\
    Months 8--10 & \textbf{Multilingual and cross-model stage:} additional-language runs, TTS vs natural stratification, secondary-model verification of core rankings. \\
    Months 10--13 & \textbf{Multimodal extensions:} interaction-value studies, modality trade-off analysis, encoder/task generalization, advanced metrics beyond faithfulness. \\
    Months 13--15 & \textbf{Other/finalization:} failure-mode closure, theory appendix, reproducibility packaging, artifact governance, and performance optimization. \\
    \bottomrule
  \end{tabularx}
\end{frame}

\begin{frame}{Risk register and mitigation}
  \begin{itemize}
    \item \textbf{Connector/model instability} \(\rightarrow\) lock fallback model candidates early.
    \item \textbf{Underpowered experiments} \(\rightarrow\) pilot-based power analysis before full runs.
    \item \textbf{Circular faithfulness evaluation} \(\rightarrow\) separate utility for SHAP computation vs evaluation.
    \item \textbf{Compute bottlenecks} \(\rightarrow\) cluster queue planning and phase-gated high-information ablations.
    \item \textbf{Cross-model claim fragility} \(\rightarrow\) explicit scoping if replication diverges.
  \end{itemize}
\end{frame}

\begin{frame}{Master's thesis and team needs}
  \begin{itemize}
    \item Separate deliverables into thesis-oriented contributions and conference-paper contributions.
    \item Add dedicated DevOps support for cluster scheduling, artifact storage, and reproducible experiment orchestration.
  \end{itemize}
\end{frame}

\begin{frame}{Conference timeline and scope fit (May 2026 -- Q1 2027)}
  \scriptsize
  \begin{tabularx}{\textwidth}{>{\raggedright\arraybackslash}p{2.7cm} >{\raggedright\arraybackslash}p{1.9cm} >{\raggedright\arraybackslash}X >{\raggedright\arraybackslash}p{1.2cm}}
    \toprule
    \textbf{Venue} & \textbf{Apply by} & \textbf{Scope link and expected deliverable} & \textbf{Fit} \\
    \midrule
    EMNLP 2026 main & May 25, 2026 & Only viable if Paper 1 results are already in hand; otherwise forces premature writing & Low (skip) \\
    BlackboxNLP 2026 & \(\sim\) mid-Aug 2026 & Primary Paper 1 target: cross-modal grounding diagnostics with archival NLP interpretability audience & High \\
    AAAI 2027 & \(\sim\) late Jul/early Aug 2026 & Secondary/main-track backup for Paper 1 variant if results mature early & Medium \\
    ICASSP 2027 & \(\sim\) early Sep 2026 & SGPA plan-B venue and optional audio-flavored Paper 1 variant (4-page format) & High \\
    NeurIPS workshops 2026 & late Aug--mid Sep 2026 & Short-form Paper 1 for additional audience exposure and feedback cycle & High \\
    ICLR 2027 & abs \(\sim\) Sep 19; full \(\sim\) Sep 24, 2026 & Primary Paper 2 benchmark submission (SV/IG/attention/LIME/ablation suite) & Medium \\
    ARR route (NAACL/EACL/ACL) & Oct 2026 / Feb 2027 cycles & Resubmission/escalation path after ICLR/workshop outcomes & Medium--High \\
    \bottomrule
  \end{tabularx}
\end{frame}

\begin{frame}{Funding and compute applications}
  \scriptsize
  \begin{tabularx}{\textwidth}{>{\raggedright\arraybackslash}p{2.9cm} >{\raggedright\arraybackslash}p{2.0cm} >{\raggedright\arraybackslash}X}
    \toprule
    \textbf{Target} & \textbf{Apply by} & \textbf{Role and gating dependency} \\
    \midrule
    PRELUDIUM 25 & Jun 16, 2026 & Feasible only with fast WUT host confirmation; otherwise skip and defer to next call \\
    OPUS 31 & Jun 16, 2026 & High-leverage path if Maria leads PI role and proposal can be aligned quickly \\
    FNP START 2027 & Q4 2026 call & Prestige track aligned with accepted/archival publication outputs \\
    LIDER UP 1/2027 & Q1 2027 (eligibility-gated) & Requires doctoral-student or WUT employment status before call opens \\
    Anthropic / OpenAI / Google credits & rolling (apply immediately) & Short-term API compute to de-risk benchmark and ablation execution \\
    PLGrid / EuroHPC access & rolling/quarterly & Primary cluster backbone for high-volume phase 2--4 experiment matrices \\
    \bottomrule
  \end{tabularx}
  \vspace{0.25em}
  \begin{itemize}
    \item \textbf{Prioritization rule:} if bandwidth is constrained, deliver Paper 1 + SGPA plan B first; defer Paper 2.
    \item \textbf{Immediate actions (next 2 weeks):} apply for rolling compute credits and PLGrid access; finalize OPUS/PRELUDIUM decision with WUT stakeholders.
  \end{itemize}
\end{frame}

\end{document}
